% Layout and typography for XeLaTeX
\usepackage{fontspec}
\usepackage{polyglossia}
\setdefaultlanguage{russian}
\setotherlanguage{english}
\IfFontExistsTF{PT Serif}{\setmainfont{PT Serif}}{\setmainfont{TeX Gyre Termes}}
\IfFontExistsTF{PT Sans}{\setsansfont{PT Sans}}{\setsansfont{TeX Gyre Heros}}
\IfFontExistsTF{PT Mono}{\setmonofont{PT Mono}}{\setmonofont{TeX Gyre Cursor}}
\usepackage[a4paper,left=30mm,right=10mm,top=20mm,bottom=20mm]{geometry}
\usepackage{setspace}
\onehalfspacing
\usepackage{indentfirst}
\setlength{\parindent}{1.25cm}
\setlength{\parskip}{0pt}
\usepackage{graphicx}
\graphicspath{{figures/}}
\usepackage{float}
\usepackage{booktabs,longtable,array,tabularx,makecell,calc}
\usepackage{caption}
\captionsetup{font=normalsize,labelfont=bf,justification=centering,singlelinecheck=false}
\usepackage{listings}
\usepackage{xcolor}
\lstdefinestyle{pythoncourse}{
  language=Python,
  basicstyle=\ttfamily\fontsize{8.5}{10}\selectfont,
  keywordstyle=\bfseries,
  commentstyle=\itshape,
  stringstyle=\ttfamily,
  numbers=left,
  numberstyle=\tiny,
  stepnumber=1,
  numbersep=7pt,
  frame=single,
  breaklines=true,
  breakatwhitespace=false,
  showstringspaces=false,
  tabsize=4,
  columns=fullflexible,
  keepspaces=true,
  xleftmargin=1em,
  framexleftmargin=0.5em,
  aboveskip=0.5\baselineskip,
  belowskip=0.5\baselineskip
}
\lstset{style=pythoncourse}
\newenvironment{listingblock}{\par\smallskip\begin{minipage}{\linewidth}}{\end{minipage}\par\smallskip}
\usepackage{titlesec}
\renewcommand{\chaptername}{Глава}
\titleformat{\chapter}[block]{\centering\bfseries\fontsize{16}{19}\selectfont}{\chaptername\ \thechapter.}{0.5em}{}
\titlespacing*{\chapter}{0pt}{-8pt}{24pt}
\titleformat{\section}[block]{\centering\bfseries\fontsize{14}{17}\selectfont}{\thesection.}{0.5em}{}
\titlespacing*{\section}{0pt}{18pt}{12pt}
\titleformat{\subsection}[block]{\bfseries\fontsize{14}{17}\selectfont}{\thesubsection.}{0.5em}{}
\titlespacing*{\subsection}{0pt}{14pt}{8pt}
\usepackage{tocloft}
\renewcommand{\contentsname}{Оглавление}
\renewcommand{\cftchapfont}{\bfseries}
\renewcommand{\cftchappagefont}{\bfseries}
\renewcommand{\cftsecfont}{\bfseries}
\renewcommand{\cftsecpagefont}{\bfseries}
\setlength{\cftbeforechapskip}{5pt}
\setcounter{secnumdepth}{2}
\setcounter{tocdepth}{2}
\usepackage{fancyhdr}
\fancyhf{}
\fancyfoot[C]{\thepage}
\renewcommand{\headrulewidth}{0pt}
\renewcommand{\footrulewidth}{0pt}
\pagestyle{fancy}
\fancypagestyle{plain}{\fancyhf{}\fancyfoot[C]{\thepage}\renewcommand{\headrulewidth}{0pt}}
\usepackage{xurl}
\usepackage{hyperref}
\hypersetup{hidelinks,pdftitle={Курсовой проект: генерация названий растений и продуктов},pdfauthor={Жаринов К. С.}}
\usepackage{microtype}
\emergencystretch=2em
